\documentclass[11pt,a4paper]{article}

% ---- Encoding & language ----
\usepackage[utf8]{inputenc}
\usepackage[T1]{fontenc}
\usepackage[english]{babel}

% ---- Layout ----
\usepackage[margin=2.5cm]{geometry}
\usepackage{parskip}
\usepackage{microtype}

% ---- Typography ----
\usepackage{mathpazo}
\linespread{1.15}

% ---- Section styling ----
\usepackage{titlesec}
\titleformat{\section}
  {\normalfont\Large\sffamily\bfseries}{\thesection}{0.7em}{}
\titleformat{\subsection}
  {\normalfont\large\sffamily\bfseries}{\thesubsection}{0.6em}{}

% ---- Math, figures, tables ----
\usepackage{amsmath, amssymb}
\usepackage{graphicx}
\usepackage{caption}
\captionsetup{font=small, labelfont={sf,bf}}
\usepackage{subcaption}
\usepackage{booktabs}
\usepackage{float}

% ---- Algorithms ----
\usepackage{algorithm}
\usepackage{algpseudocode}

% ---- References & links ----
\usepackage[numbers,square,sort&compress]{natbib}
\usepackage{xcolor}
\usepackage[colorlinks=true,
            linkcolor=black,
            citecolor=blue!45!black,
            urlcolor=blue!45!black]{hyperref}
\usepackage{enumitem}

\begin{document}
\pagenumbering{gobble}

% =============================================================================
\begin{titlepage}
\centering
\thispagestyle{empty}

{\sffamily\large NOVA Information Management School\par}
{\sffamily\normalsize Universidade Nova de Lisboa\par}

\vspace{1.5cm}
\vspace{2cm}

{\sffamily\large Text Mining\par}
{\sffamily\normalsize Academic Year 2025/26 \textbar{} Course Project\par}

\vfill

{\sffamily\Huge\bfseries Text Mining Project Final Report\par}
\vspace{0.4em}
{\sffamily\Huge\bfseries\itshape Finance Tweets Sentiment Analysis\par}
\vspace{0.4em}

\vfill

{\sffamily\large\bfseries Authors\par}
\vspace{0.4em}
\begin{tabular}{rl}
  Filipe Brand\~ao Carmo & 20240828 \\
  Marta La Feria         & 20211051 \\
  Gon\c{c}alo Bento      & 20241137 \\
  Pedro Carrasqueira     & 20250488 \\
\end{tabular}

\vspace{0.8cm}
{\sffamily Group 31\par}
\vspace{0.4em}

\vfill
\end{titlepage}

\pagenumbering{arabic}
\setcounter{page}{1}

% =============================================================================
\tableofcontents
\newpage

% =============================================================================
\section{Data Exploration}

Our exploratory data analysis of the Finance Tweets Sentiment Corpora reveals a highly compact but heavily imbalanced dataset. 
The training dataset consists of \textbf{9,543 tweets} categorized into Bearish (0), Bullish (1), and Neutral (2) sentiments.

\subsection{Class Imbalance}
\begin{itemize}[noitemsep]
    \item \textbf{Neutral}: 6,178 (64.74\%)
    \item \textbf{Bullish}: 1,923 (20.15\%)
    \item \textbf{Bearish}: 1,442 (15.11\%)
\end{itemize}

This severe majority-class bias requires strict stratified splitting and minority-class compensatory techniques (e.g., balanced class weights) to avoid deceptive accuracy metrics.

\subsection{Text Characteristics and Metadata}
Tweets are highly compressed, averaging 85.8 characters and 12.2 tokens. Financial texts are uniquely structured:
\begin{itemize}[noitemsep]
    \item \textbf{URLs}: Appear in 53.0\% of tweets, strongly indicating Neutral factual journalism (56.9\% rate).
    \item \textbf{Cashtags}: Stock tickers (e.g., \texttt{\$AAPL}) appear in 20.1\% of tweets and are strong Bullish indicators (34.1\% rate).
    \item \textbf{Non-ASCII Characters}: Over 23\% of tweets contain smart punctuation, requiring robust normalization.
\end{itemize}

% =============================================================================
\section{Corpus Split and Data Preprocessing}

To preserve label ratios, we partitioned our dataset using a strict \textbf{80/20 stratified train-validation split} (7,634 train; 1,909 validation).

We designed a robust 7-step preprocessing pipeline specifically tailored for financial microblogs:
\begin{enumerate}[noitemsep]
    \item \textbf{Unicode \& Punctuation Normalization}: Converting smart quotes and en-dashes to standard ASCII.
    \item \textbf{Lowercasing}: Collapsing identical terms.
    \item \textbf{Regex Metadata Cleaning}: Replacing sparse, unique URLs and user handles with protected placeholders (\texttt{URL\_PLACEHOLDER}, \texttt{MENTION\_PLACEHOLDER}).
    \item \textbf{NLTK TweetTokenizer}: Preserving hashtags and financial symbols.
    \item \textbf{Custom Stopwords Removal}: Removing standard stopwords plus financial noise terms (e.g., ``rt'', ``http'', ``share'').
    \item \textbf{WordNet Lemmatization}: Mapping words to roots without the destructive truncation of standard stemmers.
\end{enumerate}

% =============================================================================
\section{Feature Engineering}

We transitioned from discrete Bag-of-Words to continuous dense vector spaces, exploring multiple representations.

\subsection{Classical Vectorization}
We implemented \textbf{CountVectorizer (BoW)} and \textbf{TfidfVectorizer} for unigrams. BoW slightly outperformed TF-IDF on Macro F1 (0.6990 vs. 0.6957) because in ultra-short financial tweets, raw token presence is a stronger sentiment indicator than inverse-frequency dampened weights. 

\subsection{Word Embeddings}
We compared a \textbf{Custom Word2Vec model} (trained from scratch) against a \textbf{pre-trained GloVe-Twitter-100 model}.
\begin{itemize}[noitemsep]
    \item \textbf{OOV Rate}: Custom Word2Vec had 0\% OOV for localized tokens and placeholders, whereas GloVe missed 33\% of unique validation terms (failing on cashtags and custom jargon).
    \item \textbf{Pooling Strategies}: We explored Mean Vector Pooling vs. TF-IDF Weighted Average Pooling. Counter-intuitively, Mean Pooling outperformed TF-IDF weighting, as the latter over-amplified rare, noisy slang terms.
    \item \textbf{The Pooling Bottleneck}: Both continuous models struggled against the optimized high-dimensional TF-IDF unigram+bigram baseline because averaging word vectors destroys word-order and negation boundaries.
\end{itemize}

\noindent \textbf{[EXTRA WORK - FEATURE ENGINEERING]}: \\
To resolve static pooling limitations, we expanded our architecture to use \textbf{Dynamic Contextual Transformer Embeddings}, leveraging the self-attention mechanisms of RoBERTa and FinBERT to generate context-aware representations that natively handle negation and complex clauses.

% =============================================================================
\section{Classification Models}

We systematically benchmarked 16 classical machine learning configurations across six algorithms (KNN, Logistic Regression, MNB, MLP, Random Forest, XGBoost) using Optimized TF-IDF features.

\subsection{Prioritized Analysis: Handling Class Imbalance}
Based on the severe 65\% Neutral class imbalance, we ran an exhaustive experiment grid comparing four balancing strategies: \texttt{class\_weight='balanced'}, \texttt{oversample}, \texttt{undersample}, and \texttt{none} across model families.

\noindent \textbf{Key Findings on Imbalance Strategies (Macro F1):}
\begin{itemize}[noitemsep]
    \item \textbf{Logistic Regression (LR)}: \texttt{class\_weight} (0.7113) $>$ \texttt{oversample} (0.7085) $>$ \texttt{undersample} (0.6652) $>$ \texttt{none} (0.6417).
    \item \textbf{Multi-Layer Perceptron (MLP)}: \texttt{oversample} (0.6884) $>$ \texttt{none} (0.6830).
    \item \textbf{Random Forest (RF)}: \texttt{class\_weight} (0.6589) $>$ \texttt{oversample} (0.6326) $>$ \texttt{none} (0.3313).
    \item \textbf{Winner}: \textbf{Logistic Regression + \texttt{class\_weight='balanced'}} emerged as the absolute champion (F1 Macro = 0.7113). By mathematically scaling the loss function inversely to class frequencies, LR maximized recall on the minority Bearish (63.54\%) and Bullish (67.79\%) classes without losing Neutral accuracy. Over-sampling and Under-sampling produced competitive but slightly inferior bounds due to noise replication and data loss, respectively.
\end{itemize}

\subsection{Contextual Transformer Encoders}
Going beyond classical ML, we implemented and fine-tuned bi-directional Transformer Encoder models to overcome the limitations of static word embeddings. These mandatory architectural variations included:
\begin{enumerate}[noitemsep]
    \item \textbf{DistilBERT (\texttt{distilbert-base-uncased})}: A lightweight and highly efficient transformer. It proved to be the optimal architecture for our compact dataset, converging with high stability and achieving the absolute highest overall performance with a Macro F1 of \textbf{0.8236}.
    \item \textbf{Twitter-RoBERTa-base-sentiment}: Pre-trained on 58M tweets, this model natively handles informal syntax, slang, and emoticons. It successfully generalized to our corpus, achieving a strong Macro F1 of \textbf{0.8011}.
\end{enumerate}

\vspace{0.8em}
\noindent \textbf{[EXTRA WORK - ADVANCED TRANSFORMER ENCODERS]}:\\
To push the boundaries of domain adaptation and complex architectures, we explicitly implemented and fine-tuned two additional models:
\begin{enumerate}[noitemsep]
    \item \textbf{FinBERT (ProsusAI)}: Pre-trained on formal financial texts, excelling in precise financial jargon extraction. However, it struggled with the informal grammatical nature of microblogs, dropping to an F1 of 0.5911.
    \item \textbf{DeBERTa-v3-base}: An advanced model utilizing disentangled attention. Interestingly, it collapsed during fine-tuning (F1 = 0.0920), indicating a high sensitivity to hyperparameters on small datasets.
\end{enumerate}
For all external models, we programmatically handled mismatched label schemas by safely re-initializing their classification heads, performing a full end-to-end fine-tuning of all network weights to maximize domain adaptation.

\subsection{Encoder as Feature Extractor (Frozen Embeddings)}
To satisfy the architectural variation criteria, we also implemented a hybrid ``Encoder-as-Features'' pipeline. In this setup, we completely froze the pre-trained transformer backbone (preventing any gradient updates) and extracted the contextualized embeddings as static features to feed into a downstream classical classifier (e.g., Logistic Regression). While this approach drastically reduced computational overhead and training time compared to full fine-tuning, the constrained representational power produced sub-optimal F1 scores. This explicitly confirmed that unfreezing the transformer weights is strictly necessary for deep financial domain adaptation.

\vspace{0.8em}
\noindent \textbf{[EXTRA WORK - DECODER-BASED CLASSIFICATION (QWEN)]}:\\
Expanding our exploration into cutting-edge Generative AI, we designed a pure sequence-to-sequence inference pipeline using a Large Language Model. We deployed \textbf{Qwen2.5-1.5B-Instruct} strictly as a zero-gradient decoder using a meticulously engineered \textbf{stratified 6-shot chat template} (2 representative examples per class). 
By configuring the model for deterministic greedy generation (\texttt{max\_new\_tokens=4}, \texttt{dtype=bf16}) and strictly enforcing the label schema, we achieved a perfect \textbf{0.00\% parse failure rate} during batch inference. 
Although its Macro F1-Score (0.5592) trailed fully fine-tuned bidirectional encoders---highlighting the extreme difficulty of few-shot prompt-based classification on highly ambiguous short financial texts---this robust implementation successfully demonstrated the viability of utilizing foundational decoders via prompt engineering without the need for computationally expensive backpropagation.

% =============================================================================
\section{Evaluation and Results}

Table \ref{tab:model_comparison} presents the comprehensive evaluation metrics across the selected baseline and advanced architectures. All models are evaluated across four mandatory metrics: Accuracy, Macro Precision, Macro Recall, and Macro F1-Score.

\begin{table}[H]
    \centering
    \begin{tabular}{lcccc}
        \toprule
        \textbf{Model} & \textbf{Accuracy} & \textbf{Precision} & \textbf{Recall} & \textbf{F1-Score} \\
        \midrule
        Logistic Regression (Classical Baseline) & 0.7810 & 0.7041 & 0.7201 & 0.7113 \\
        \textbf{DistilBERT (Champion)} & \textbf{0.8628} & \textbf{0.8361} & 0.7964 & \textbf{0.8236} \\
        Twitter-RoBERTa & 0.8400 & 0.7968 & \textbf{0.8101} & 0.8011 \\
        FinBERT & 0.7400 & 0.5934 & 0.5915 & 0.5911 \\
        Qwen2.5-1.5B (Decoder-based) & 0.6350 & 0.6228 & 0.6046 & 0.5592 \\
        DeBERTa-v3-base & 0.1600 & 0.0533 & 0.3333 & 0.0920 \\
        \bottomrule
    \end{tabular}
    \caption{Comparative Performance Metrics of Implemented Models.}
    \label{tab:model_comparison}
\end{table}

\subsection{DistilBERT Confusion Matrix}
To further visualize the predictions of our highest performing model across the Bearish, Bullish, and Neutral classes, Figure \ref{fig:confusion_matrix} displays the confusion matrix for the fine-tuned DistilBERT.

\begin{figure}[H]
    \centering
    \includegraphics[width=0.65\textwidth]{../outputs/confusion_matrix.png}
    \caption{Confusion Matrix for the Champion Model (DistilBERT).}
    \label{fig:confusion_matrix}
\end{figure}

\subsection{Error Analysis \& Failure Modes}
Despite high performance, confusion matrix analysis reveals distinct limitations of linear BoW models:
\begin{enumerate}[noitemsep]
    \item \textbf{Contextual Negation Blindness}: Failing to resolve sequential dependencies (e.g., ``not to like'', ``fail to halt''), leading to polarity inversions.
    \item \textbf{Factual Reporting vs. Sentiment}: The model struggles with boundary ambiguity, aggressively flagging factual positive corporate news (often labeled Neutral) as Bullish.
    \item \textbf{Clause Overriding}: In complex tweets (e.g., \textit{``profit rise fails to salvage''}), the sum of positive token weights overrides single critical negative tokens.
\end{enumerate}
These failures definitively justify our advanced transition into context-aware deep learning models (FinBERT/RoBERTa) capable of gating semantic flow.

% =============================================================================
\section{Final Model Selection}

Based on the empirical evidence from our extensive experiments, the final champion model selected for our official pipeline is the \textbf{Fine-Tuned DistilBERT} (\texttt{distilbert-base-uncased}).

\subsection{Justification:}

\begin{enumerate}[noitemsep]
    \item \textbf{Absolute Performance Winner}: DistilBERT achieved the highest overall Macro F1-Score (\textbf{0.8236}) across all 30+ tested configurations, dramatically outperforming the best classical machine learning model (Logistic Regression, F1=0.7113).
    \item \textbf{Contextual Semantic Superiority}: Our error analysis revealed that classical TF-IDF pipelines suffer from ``Contextual Negation Blindness'' and clause overriding because they ignore word order. DistilBERT's bi-directional self-attention layers natively resolve these syntactical dependencies (e.g., cleanly distinguishing \textit{``fails to drop''} from \textit{``drops''}), solving the primary failure mode of our baseline models.
    \item \textbf{Imbalance Resilience}: DistilBERT attained strong recall on the difficult minority Bearish class (Class F1 = 0.734) simply through superior dense representation learning, entirely avoiding the need for artificial SMOTE oversampling or mathematical loss-scaling hacks (\texttt{class\_weight='balanced'}).
    \item \textbf{Efficiency vs. Domain Specificity}: Counter-intuitively, DistilBERT out-performed heavily domain-specific financial models like FinBERT (0.5911) and massive social media models like Twitter-RoBERTa (0.8011). DistilBERT proved to be the optimal ``Goldilocks'' model: it has enough general English representation to handle syntax, but is lightweight enough to achieve fast, highly stable convergence during fine-tuning on our compact dataset without catastrophically overfitting or collapsing (like DeBERTa-v3).
\end{enumerate}

Consequently, \texttt{distilbert-base-uncased} has been fully integrated into our Agentic Dispatcher (\texttt{autotune.py}) and is deployed as the final model for \texttt{tm\_final\_xx.ipynb}.

\end{document}
